\clearpage
\section{Geometric-Median Prototypes and the Weiszfeld Solver}
\label{sec:kdgrp-weiszfeld-appendix}

This appendix makes the robust-prototype calculation independently
understandable. Mathematical identities are separated from implementation
choices. The latter must be copied from the frozen production audit before the
report is submitted.

\subsection{Arithmetic mean versus geometric median}

For calibration features $\{\bm f_n\}_{n=1}^{N}$, the arithmetic mean minimizes
the squared-distance objective
\begin{equation}
    J_2(\bm c)=\sum_{n=1}^{N}\|\bm f_n-\bm c\|_2^2.
\end{equation}
Differentiating gives
\begin{equation}
    \nabla J_2(\bm c)
    =2N\bm c-2\sum_{n=1}^{N}\bm f_n.
\end{equation}
Setting the gradient to zero yields the explicit solution
\begin{equation}
    \bar{\bm f}=\frac{1}{N}\sum_{n=1}^{N}\bm f_n.
\end{equation}
Because the objective squares each residual, the magnitude of one point's
gradient contribution grows linearly with its distance from the current centre.

The geometric median instead minimizes
\begin{equation}
    J_1(\bm c)=\sum_{n=1}^{N}\|\bm f_n-\bm c\|_2.
\end{equation}
For $\bm c\neq\bm f_n$, each contribution satisfies
\begin{equation}
    \left\|
    \frac{\bm c-\bm f_n}{\|\bm c-\bm f_n\|_2}
    \right\|_2=1.
\end{equation}
Thus, an extreme point is not ignored, but its influence does not grow merely
because it is farther from the centre. This property motivates its use as a
robust reference for a calibration population containing rare spatial-feature
states.

\subsection{First-order condition and the implicit solution}

When the candidate centre does not coincide with an observation,
\begin{equation}
    \nabla J_1(\bm c)
    =\sum_{n=1}^{N}
      \frac{\bm c-\bm f_n}{\|\bm c-\bm f_n\|_2}.
\end{equation}
At a differentiable optimum $\bm c^\star$,
\begin{equation}
    \sum_n
    \frac{\bm c^\star-\bm f_n}
         {\|\bm c^\star-\bm f_n\|_2}=0.
\end{equation}
Moving the feature terms to the right and collecting $\bm c^\star$ gives
\begin{equation}
    \bm c^\star
    \sum_n\frac{1}{\|\bm c^\star-\bm f_n\|_2}
    =
    \sum_n\frac{\bm f_n}{\|\bm c^\star-\bm f_n\|_2},
\end{equation}
and hence
\begin{equation}
    \boxed{
    \bm c^\star=
    \frac{\displaystyle\sum_n
          \bm f_n/\|\bm c^\star-\bm f_n\|_2}
         {\displaystyle\sum_n
          1/\|\bm c^\star-\bm f_n\|_2}.}
\end{equation}
The unknown centre occurs inside every denominator, so this is a nonlinear
fixed-point equation rather than an explicit solution. A generic
multidimensional geometric median has no general arithmetic-mean-like formula
and is normally solved iteratively. This statement allows the familiar
one-dimensional median and specially symmetric configurations as exceptions.

\subsection{Weiszfeld fixed-point and IRLS derivation}

Weiszfeld iteration replaces the unknown weights with weights computed from the
current iterate:
\begin{equation}
    w_n^{(t)}=
    \frac{1}{\|\bm f_n-\bm c^{(t)}\|_2+\varepsilon},
    \qquad
    \bm c^{(t+1)}=
    \frac{\sum_nw_n^{(t)}\bm f_n}{\sum_nw_n^{(t)}}.
\end{equation}
The displayed $+\varepsilon$ is the currently documented stabilized form; the
final report must match the exact frozen implementation.

The same update can be understood as iteratively reweighted least squares
(IRLS). Let
\begin{equation}
    r_n^{(t)}=\|\bm f_n-\bm c^{(t)}\|_2.
\end{equation}
For $r_n^{(t)}>0$, the tangent upper bound for the square root gives
\begin{equation}
    \|\bm f_n-\bm c\|_2
    \leq
    \frac{\|\bm f_n-\bm c\|_2^2}{2r_n^{(t)}}
    +\frac{r_n^{(t)}}{2}.
\end{equation}
Summing these bounds produces the quadratic surrogate
\begin{equation}
    Q(\bm c\mid\bm c^{(t)})
    =\sum_n
      \frac{\|\bm f_n-\bm c\|_2^2}{2r_n^{(t)}}
      +\text{constant}.
\end{equation}
Minimizing this weighted least-squares problem yields
\begin{equation}
    \bm c^{(t+1)}=
    \frac{\sum_n\bm f_n/r_n^{(t)}}
         {\sum_n1/r_n^{(t)}}.
\end{equation}
The stabilized implementation prevents a singular inverse distance when an
iterate approaches an observation.

\subsection{Production initialization}

\evidenceplaceholder{KDGRP-IMPLEMENTATION-02: exact initializer}{
Insert the implementation-backed formula for $\bm c^{(0)}$. State whether it
is the arithmetic mean, coordinate-wise median, a selected calibration vector,
or another deterministic procedure. Explain that it initializes the solver and
does not redefine the geometric-median objective. Record any seed and the
check used when the initializer coincides with a calibration observation.}

\subsection{Numerical stabilization and stopping}

\evidenceplaceholder{KDGRP-IMPLEMENTATION-03: stabilization and convergence}{
Insert the exact inverse-distance convention, numerical $\varepsilon$,
convergence tolerance, maximum inner iterations, stopping norm, dtype/device,
and near-coincidence handling from the frozen production implementation. If the
code uses $1/\max(r,\varepsilon)$ rather than $1/(r+\varepsilon)$, replace the
displayed equations consistently.}

\subsection{Extension to multiple modes}

For $K>1$, the frozen method alternates between nearest-centre assignment,
\begin{equation}
    a_n=\arg\min_k\|\bm f_n-\bm c_k\|_2,
\end{equation}
and a geometric-median update within each assigned cluster,
\begin{equation}
    \bm c_k=\arg\min_{\bm c}
      \sum_{n:a_n=k}\|\bm f_n-\bm c\|_2.
\end{equation}
The joint assignment-and-centre objective is not globally convex. The
production initialization, tie handling, empty-cluster policy, and stopping
rule therefore matter for reproducibility.

\evidenceplaceholder{KDGRP-IMPLEMENTATION-04: exact multi-mode solver}{
Insert the audited initial $K$-centre algorithm, seed, assignment tie rule,
empty-cluster handling, per-cluster Weiszfeld initializer, outer stopping rule,
maximum outer iterations, restart policy, deterministic ordering requirements,
and whether lower-$K$ solutions warm-start higher $K$. Do not claim global
optimality.}

\subsection{Connection to the explanation score}

The solver is used only during label-free calibration to estimate recurrent
reference modes. Once these modes are frozen, explanation inference uses
\begin{equation}
    R_K(p)=\min_k\|\bm f_p-\bm c_k\|_2.
\end{equation}
No iterative optimization is required for an individual image. Inference needs
only feature extraction, distances to the frozen centres, nearest-mode
selection, calibration through the frozen CDF, and spatial reconstruction.
